% §II Related Work (~0.75 page). 4 clusters per OUTLINE §5-II. Differentiation via 토론 1
% (3-layer claim discipline: their evidence doesn't establish causal use; ours is an existence
% proof of bypass; falsifiable prediction, not accusation).
% 08-11 reviewer-proofing (reference/ANALYSIS.md 전수조사 기반): near-miss 구분 추가 —
% TA-VLA(노이즈는 baseline 토큰, HSIC는 상관), Tactile-VLA(dose-response는 language 입력),
% ForceFlow(마스킹은 vision 입력), ForceSight(개입은 출력 force goal, readout은 SR).
% hedge "rollout force statistics" §I과 미러링. FuSe/FM-VLA = 분야가 인지하고도 SR로만 검증한 사례.
\section{Related Work}
\label{sec:related}

\textbf{Force- and tactile-conditioned VLA policies.}
ForceVLA fuses wrist force as dedicated tokens through a mixture-of-experts
backbone~\cite{forcevla}; TA-VLA injects torque adapters into the action
decoder~\cite{tavla}; PhaForce schedules policy phases on contact and force
signals~\cite{phaforce}; ForceFlow~\cite{forceflow}, FARM~\cite{farm}, and
FILIC~\cite{filic} condition diffusion or impedance-loop policies on force and tactile
input; TacForeSight learns force-conditioned tactile prediction~\cite{tacforesight}.
Each commits to a way of getting force into the network---new tokens, new experts, new
adapters, a new decoder---and each is validated by task success and retraining ablations,
sometimes alongside rollout force statistics~\cite{phaforce,forceflow}. Our pathway is the
lightest of these choices: a computed condition modulating the state token of an otherwise
unchanged backbone, attachable after training. We also evaluate differently: rather than
comparing differently trained policies by success, we intervene on the force input of a
trained policy and read the action, which exposes the \emph{form} of the response---whether
it grows or fades with force---that success rates average over. The closest such analyses
stop short of the force input: TA-VLA perturbs input tokens of its baseline~\cite{tavla},
Tactile-VLA's dose-response intervenes on the instruction~\cite{tactilevla}, and ForceFlow
masks its trained policy's visual input~\cite{forceflow}.

\textbf{Intermediate condition representations.}
Injecting a compact, task-level representation between observation and action is a recurring
design: RT-Affordance conditions on predicted affordance plans~\cite{rtaffordance}, and FiLM
modulation~\cite{perez2018film} is the standard mechanism for conditioning a backbone on side
information. Our \chat{} shares the shape but differs in role: it is computed from signals the
policy already receives rather than learned, so it adds no information and its meaning is fixed
before training.

\textbf{Causal confusion and shortcut learning.}
That imitators latch onto spurious correlates is well established: causal confusion in
IL~\cite{dehaan2019causal}, copycat behavior on observation histories~\cite{wen2020copycat},
and shortcut learning in deep networks broadly~\cite{geirhos2020shortcut,hermann2023foundations}.
Recent work documents shortcuts in generalist robot policies~\cite{shortcutgeneralist} and uses
targeted interventions to expose or correct them~\cite{interventionil2025,interventionil2023}.
The symptoms are on record in force-FiLM policies themselves: FuSe observes fine-tuned
generalists ignoring newly added sensors~\cite{fuse}, and FM-VLA diagnoses a history-length
shortcut in its own force stream and patches it with augmentation---verifying the fix, again,
only by task success~\cite{fmvla}.

\textbf{Privileged supervision and interventional evaluation.}
Privileged-teacher distillation~\cite{chen2020lbc,privilegedaction} assumes access to a signal
worth distilling; our probes instead ask how a deployed policy responds to the signal it already
has. The closest interventional precedent is ForceSight, which ablates predicted force
\emph{goals} at its controller under matched initial conditions~\cite{forcesight}; the
intervention sits on the output side of a model that takes no force input, and reads out
end-of-episode success rather than the action. Our press-retreat demonstrations are a data-side
decorrelation in the spirit of fixes for causal confusion~\cite{dehaan2019causal}, motivated
by a physical argument (Sec.~\ref{sec:discussion}).
